% !TEX program = xelatex
% Статья для сборника НИТУ «МИСиС» (шаблон: sadekovlar/misis-article-latex)


\documentclass{misisarticle}

\usepackage{tikz}
\usetikzlibrary{arrows.meta, positioning}


\grnti{28.23, 06.81}
\asjc{1702, 1406}
\oecd{1.02, 5.02}

\title{Персонализация маркетинговых коммуникаций в экосистеме сервисов на основе контекстных бандитов и больших языковых моделей}


\misisauthor{В.~М.~Гришин}{кафедра инженерной кибернетики НИТУ «МИСиС»}{Москва, Россия}{student@edu.misis.ru}

\abstract{%
  Рассматривается задача персонализации маркетинговых коммуникаций в приложении с несколькими сервисами: для каждого клиента нужно выбрать сервис и предложение, канал, время, содержание сообщения или отказаться от контакта. Распространённые правила и модели отклика не учитывают, совершил бы клиент целевое действие без сообщения, и не согласуют контакты разных сервисов между собой. Цель работы~--- увеличить дополнительный эффект коммуникаций при ограниченном числе контактов. Предлагается модульная политика: контекстный бандит с сэмплированием Томпсона выбирает компоненты сообщения с учётом молчания как отдельного действия и утомления клиента, последовательность контактов планируется на скользящем горизонте, а большая языковая модель формирует кандидатов последовательности и текст в рамках выбранных компонентов, но не принимает итоговое решение. Проверка описана для полусинтетической среды: отклик калибруется по открытому набору Open Bandit Dataset, а эффекты каналов, времени и частоты задаются симулятором с известной истиной. Политика сравнивается с базовыми стратегиями по накопленному сожалению и точности оценки эффекта, тексты~--- по доле утверждений, подтверждённых исходными данными.
}

\keywords{%
  контекстные бандиты, сэмплирование Томпсона, модульное пространство действий, инкрементальный эффект, маркетинговые коммуникации, большие языковые модели
}

\begin{document}

\maketitle

% ============================================================
\section{Введение}

Приложение, объединяющее несколько сервисов, может обращаться к клиенту через пуш-уведомления, электронную почту и внутренние поверхности. Обычно такими рассылками управляют маркетологи вручную: они задают сегменты, расписания и шаблоны. При этом содержание, время, частота и текст сообщения слабо персонализированы, а кампании разных сервисов конкурируют за внимание одного клиента~\cite{abboud2025}.

У автоматизации выбора есть три трудности. Во-первых, высокий отклик не означает пользу: клиент мог совершить покупку и без сообщения, поэтому нужно оценивать дополнительный, или инкрементальный, эффект~\cite{sato2019}. Во-вторых, решения связаны между собой. Удачное сегодня сообщение может снизить отклик на следующие, поэтому выбор отдельного лучшего контакта не гарантирует лучшую последовательность~\cite{yancey2020}. В-третьих, большие языковые модели (LLM) упрощают написание текста, но их выводы о клиенте могут не подтверждаться его данными~\cite{fostiropoulos2026, zhang2026}.

Настоящая работа посвящена подходу, который разделяет статистический выбор действия и генерацию. Планируемый вклад работы состоит в следующем:
\begin{itemize}
  \item предлагается модульная политика, которая согласованно выбирает сервис, канал, время и смысловые компоненты сообщения, считая молчание отдельным действием и оптимизируя инкрементальный эффект за вычетом стоимости контакта;
  \item предлагается способ подключения LLM, при котором модель предлагает кандидатов последовательности и пишет текст, а итоговый выбор и проверка остаются за байесовской моделью и детерминированными правилами;
  \item описывается полусинтетическая среда с известным истинным эффектом и методика сравнения политики с базовыми стратегиями.
\end{itemize}

В статье изложены анализ работ, постановка задачи и описание методов. Эксперименты на момент написания ещё не проведены, поэтому в разделе~\ref{sec:methods} приведена методика их проведения, а численные результаты не приводятся.

% ============================================================
\section{Анализ работ}

\subsection{Контекстные бандиты и сэмплирование Томпсона}

Контекстный бандит на каждом шаге наблюдает контекст, выбирает действие и получает награду только за него, поэтому ему приходится балансировать между исследованием и использованием. В работе~\cite{li2010} алгоритм LinUCB применён к рекомендации новостей, а в~\cite{li2011} предложена несмещённая офлайн-оценка таких алгоритмов по логам со случайным выбором. Сэмплирование Томпсона~\cite{thompson1933} выбирает действие согласно выборке из апостериорного распределения параметров; его свойства для контекстной постановки с линейной наградой изучены в~\cite{agrawal2013}, а практические вопросы разобраны в обзоре~\cite{russo2018}.
\subsection{Составные действия и последовательности}

Сообщение состоит из нескольких независимо выбираемых компонентов, и число их сочетаний быстро растёт. В~\cite{hill2017} для многомерной оптимизации страницы использованы сэмплирование Томпсона, модель с парными взаимодействиями и приближённый поиск лучшего сочетания. В~\cite{mcinerney2018} бандит решает совместную задачу выбора рекомендаций и их объяснений. Повторяемость контактов учитывает~\cite{yancey2020}: в задаче регулярных уведомлений часть действий временно недоступна, а эффект повторно использованного шаблона восстанавливается со временем.

\subsection{Оценка инкрементального эффекта}

Метаалгоритмы оценки неоднородного эффекта воздействия описаны в~\cite{kunzel2019}. В~\cite{sato2019} показано, как оценивать и оптимизировать рекомендации по приросту, а не по отклику. Наиболее близкая к нашей постановке работа~\cite{abboud2025} формулирует оркестрацию кросс-канальных коммуникаций как последовательное принятие решений: модульная политика максимизирует инкрементальное вовлечение, индивидуальный эффект оценивается через схему «разность разностей», а исследование обеспечивается сэмплированием Томпсона. Подход внедрён на 150~млн пользователей. Вместе с тем планирование согласованной последовательности на заданный горизонт и генерация текста с помощью LLM не являются в этой работе основным предметом.

Для офлайн-оценки политик по логам применяются взвешивание по обратной вероятности и двойная устойчивость~\cite{dudik2011}. Публичный набор Open Bandit Dataset~\cite{saito2021} содержит логи нескольких политик на одной платформе и позволяет сравнивать такие оценки.

\subsection{Большие языковые модели в персонализации}

Эксперименты~\cite{matz2024} показывают, что персонализированные генеративной моделью сообщения могут быть убедительнее неперсонализированных. Гибрид~\cite{song2025} ближе всего к предлагаемому разделению ролей: бандит выбирает тип воздействия, а LLM пишет текст внутри него. Полевое сравнение~\cite{hofer2026} в задаче мотивации к физической активности показало, что сообщения LLM оцениваются выше шаблонов, однако бандит не добавил воспринимаемой полезности; при этом LLM без ограничений концентрировалась на одной технике, а бандит обеспечивал систематическое исследование. Эти две работы измеряют прежде всего субъективные оценки, а не дополнительный эффект на целевое действие.

Надёжность LLM в понимании пользователя ограничена. Бенчмарк GISTBench~\cite{fostiropoulos2026} отдельно оценивает, подтверждаются ли предсказанные интересы историей взаимодействий. В~\cite{zhang2026} показано, что на данных реальных людей модели с трудом извлекают атрибуты пользователя, а персонализированные ответы люди оценивают не выше обычных. Поэтому мы не поручаем LLM выбирать действие и проверяем её утверждения по исходным данным.

\subsection{Связанные работы сборника}

Из работ сборника кафедры ближе всего к нашей задаче~\cite{sosedka2024}: там LLM выполняет семантический анализ сообщений Telegram-каналов, а на полученных данных строится классификатор для рекомендательной системы. В~\cite{chetvertkov2026} сравниваются мультимодальные архитектуры текстового поиска 3D-моделей LEGO на наборе данных с описаниями, сгенерированными LLM. В~\cite{zaitsev2026} при восстановлении сцены по одному изображению среди прочих подходов рассмотрен конвейер на основе больших языковых моделей. Предметная область этих работ иная, однако близка их схема: сравнение нескольких подходов по единой методике. Устойчивое решение задачи квадратичной минимизации с неравномерно возмущённым оператором рассмотрено в~\cite{artemyeva2022}; близкая ситуация возникает при обновлении модели отклика по зашумлённым данным.

\subsection{Выводы из анализа}

В табл.~\ref{tab:related} работы сопоставлены с предлагаемым подходом. В просмотренной литературе не найдено работы, которая одновременно оптимизировала бы инкрементальный эффект модульного выбора, планировала последовательность на горизонт и использовала LLM под контролем статистической модели. Этот вывод основан на ограниченном поиске и будет уточняться.

\begin{table}[htbp]
  \caption{Сопоставление с близкими работами}
  \label{tab:related}
  \centering
  \small
  \begin{tabular}{>{\raggedright\arraybackslash}p{0.34\columnwidth}cccc}
    \toprule
    \textbf{Работа} & \textbf{Мод.} & \textbf{Инкр.} & \textbf{Гор.} & \textbf{LLM} \\
    \midrule
    LinUCB~\cite{li2010} & --- & --- & --- & --- \\
    Многомерный бандит~\cite{hill2017} & + & --- & --- & --- \\
    Восстановление~\cite{yancey2020} & --- & --- & --- & --- \\
    Agentic~\cite{abboud2025} & + & + & --- & --- \\
    cMABxLLM~\cite{song2025} & --- & --- & --- & + \\
    Предлагаемый подход & + & + & + & + \\
    \bottomrule
  \end{tabular}

  \vspace{1mm}
  \footnotesize Мод.~--- модульное действие; Инкр.~--- целевой инкрементальный эффект; Гор.~--- планирование на горизонт; LLM~--- генерация текста.
\end{table}

% ============================================================
\section{Постановка задачи}

Для клиента в день $t = 1, \dots, T$ известен контекст $x_t$: профиль, история заказов и контактов, сведения о сервисах и бизнес-ограничения. Действие $a = (p, c, s, \ell)$ составлено из модулей: $p$~--- сервис и предложение, $c$~--- канал, $s$~--- временной слот, $\ell$~--- смысловые метки сообщения (например, тон и стимул). Отдельное действие $\varnothing$ означает молчание. Множество $\mathcal{A}(x_t)$ содержит только допустимые действия: есть согласие на канал, предложение доступно, чувствительные категории исключены.

Обозначим через $Y_t(a) \in \{0, 1\}$ потенциальный исход~--- совершение целевого действия в окне $W$ дней после решения. Инкрементальный эффект равен
\begin{equation}
  \tau(x, a) = \mathbb{E}\left[\,Y(a) - Y(\varnothing) \mid x\,\right], \quad \tau(x, \varnothing) = 0.
  \label{eq:effect}
\end{equation}
Награда учитывает ценность целевого действия $v_p$ и стоимость контакта $\kappa$:
\begin{equation}
  R(x, a) = v_p\,\tau(x, a) - \kappa(x, a), \quad \kappa(x, \varnothing) = 0.
  \label{eq:reward}
\end{equation}
Требуется найти политику $\pi$, максимизирующую суммарную награду
\begin{equation}
  \max_{\pi}\ \mathbb{E}_{\pi} \sum_{t=1}^{T} R(x_t, a_t)
  \label{eq:objective}
\end{equation}
при условиях $a_t \in \mathcal{A}(x_t)$ и не более $F$ контактов за любые семь дней подряд. Контекст зависит от прошлых действий, поэтому жадный выбор на каждом шаге в общем случае неоптимален.

Качество обучения измеряется накопленным сожалением относительно оракула $\pi^{*}$, которому известны истинные параметры среды:
\begin{equation}
  \begin{aligned}
    \mathrm{Reg}(T) ={}& \mathbb{E}_{\pi^{*}} \sum_{t=1}^{T} R(x_t, a_t) \\
    &- \mathbb{E}_{\pi} \sum_{t=1}^{T} R(x_t, a_t).
  \end{aligned}
  \label{eq:regret}
\end{equation}

К тексту предъявляется отдельное требование: каждое его фактическое утверждение должно подтверждаться описанием предложения или данными клиента. Эффект конкретной формулировки в симуляции не определён, поэтому качество текста оценивается отдельно от награды~\eqref{eq:reward}.

Проверяются четыре гипотезы:
\begin{itemize}
  \item H1: модульная политика достигает меньшего сожаления, чем бандит с независимой рукой на каждое сочетание;
  \item H2: оптимизация эффекта~\eqref{eq:effect} с молчанием даёт большую истинную награду при меньшем числе контактов, чем оптимизация отклика;
  \item H3: планирование на горизонт превосходит жадный выбор при заметном утомлении клиента;
  \item H4: кандидаты от LLM не хуже кандидатов лучевого поиска при той же оценке, а проверка снижает долю неподтверждённых утверждений в тексте.
\end{itemize}

% ============================================================
\section{Описание методов}
\label{sec:methods}

Общая схема показана на рис.~\ref{fig:pipeline}. Значения параметров являются исходными и будут уточняться в пробных запусках.

\begin{figure}[htbp]
  \centering
  \begin{tikzpicture}[
      font=\footnotesize,
      box/.style={draw, rounded corners=1.5pt, align=center, text width=0.86\columnwidth, inner sep=3pt},
      arr/.style={-{Latex}, thick}]
    \node[box] (a) {Контекст $x_t$: профиль, история, допустимые действия $\mathcal{A}(x_t)$};
    \node[box, below=4mm of a] (b) {Выборка параметров $\hat{w}$ из апостериорного распределения};
    \node[box, below=4mm of b] (c) {Кандидаты последовательностей: лучевой поиск и LLM};
    \node[box, below=4mm of c] (d) {Оценка кандидатов по $\hat{w}$ и выбор первого действия};
    \node[box, below=4mm of d] (e) {Текст от LLM в рамках меток; проверка, иначе шаблон};
    \node[box, below=4mm of e] (f) {Отправка или молчание; журнал действия и вероятности};
    \node[box, below=4mm of f] (g) {Исход через $W$ дней; обновление апостериорного распределения};
    \draw[arr] (a) -- (b);
    \draw[arr] (b) -- (c);
    \draw[arr] (c) -- (d);
    \draw[arr] (d) -- (e);
    \draw[arr] (e) -- (f);
    \draw[arr] (f) -- (g);
  \end{tikzpicture}
  \caption{Цикл принятия решения по клиенту}
  \label{fig:pipeline}
\end{figure}

\subsection{Модель отклика и оценка эффекта}

Вместо отдельной оценки для каждого сочетания используется байесовская логистическая регрессия по компонентам действия. Для $a \ne \varnothing$
\begin{equation}
  \begin{aligned}
    \eta_w(x, a) ={}& w_0^{\top} x + \sum_{m} w_m^{\top} \phi_m(x, a_m) \\
    &+ \sum_{(m, n) \in \mathcal{I}} w_{mn}^{\top} \phi_{mn}(a_m, a_n) + \gamma^{\top} g(h),
  \end{aligned}
  \label{eq:model}
\end{equation}
а для молчания $\eta_w(x, \varnothing) = w_0^{\top} x$. Здесь $\phi_m$~--- признаки модуля $m$ и их произведения с контекстом, $\mathcal{I}$~--- заранее выбранные пары модулей. Признаки утомления $g(h)$ строятся по истории $h$: число контактов за семь дней и логарифм времени с последнего использования той же метки, что отражает восстановление эффекта~\cite{yancey2020}. Вероятность исхода равна $\sigma(\eta_w)$, а оценка эффекта
\begin{equation}
  \tau_w(x, a) = \sigma\!\left(\eta_w(x, a)\right) - \sigma\!\left(\eta_w(x, \varnothing)\right).
  \label{eq:tauhat}
\end{equation}

Базовая вероятность без контакта оценивается по периодам молчания. Чтобы такие данные были и у активных клиентов, с вероятностью $\varepsilon_0$ решение политики заменяется на $\varnothing$ независимо от контекста. Такая случайная контрольная группа уменьшает число контактов, зато эффект опирается не только на допущения модели.

Рассматриваются общая для всех клиентов модель и иерархическая модель с личными отклонениями $w_u = w + \delta_u$, $\delta_u \sim \mathcal{N}(0, \nu^2 I)$. При нескольких контактах в месяц личный слой может давать больше шума, чем пользы, поэтому порог объёма данных, при котором он окупается, определяется экспериментально. Аналогично проверяется, не сходится ли аддитивная модель без взаимодействий к неверному сочетанию.

\subsection{Обучение и выбор действия}

Апостериорное распределение приближается методом Лапласа $w \sim \mathcal{N}(\mu, \Sigma)$. После наступления окна $W$ для пакета исходов параметр $\mu$ уточняется шагами Ньютона, а точность~--- по формуле
\begin{equation}
  \Sigma_{\mathrm{new}}^{-1} = \delta\,\Sigma_{\mathrm{old}}^{-1} + \sum_{i} \sigma_i (1 - \sigma_i)\, z_i z_i^{\top},
  \label{eq:update}
\end{equation}
где $z_i$~--- вектор признаков наблюдения, $\sigma_i$~--- вероятность исхода в текущей точке, а множитель $\delta \le 1$ вводит забывание для нестационарных предпочтений.

В сэмплировании Томпсона~\cite{thompson1933, russo2018} для клиента берётся выборка $\hat{w} \sim \mathcal{N}(\mu, \Sigma)$, и действия оцениваются как
\begin{equation}
  \hat{R}(x, a) = v_p\,\tau_{\hat{w}}(x, a) - \kappa(x, a).
  \label{eq:rhat}
\end{equation}
Молчание имеет оценку ноль, поэтому контакт отправляется, только если хотя бы одна выборочная оценка положительна. Лучшее сочетание находится перебором допустимых действий, а при большом их числе~--- покоординатным подъёмом, как в~\cite{hill2017}. Вероятности выбора действий оцениваются методом Монте-Карло и записываются в журнал для последующей офлайн-оценки.

\subsection{Планирование последовательности}

Жадный выбор~\eqref{eq:rhat} не учитывает, что контакт сегодня снизит отклик завтра. Поэтому последовательности $\xi = (a_t, \dots, a_{t+H-1})$ оцениваются при той же выборке $\hat{w}$:
\begin{equation}
  J_{\hat{w}}(\xi) = \sum_{k=0}^{H-1} \lambda^{k}\, \hat{R}\left(x_{t+k}, a_{t+k}\right),
  \label{eq:plan}
\end{equation}
где $\lambda \le 1$~--- коэффициент дисконтирования, а признаки утомления в $x_{t+k}$ пересчитываются с учётом уже включённых в $\xi$ действий. Исходное значение горизонта $H = 7$ дней. Выполняется только первое действие; на следующий день план строится заново, поскольку поступают новые данные.

Кандидаты поступают из двух источников: лучевого поиска и LLM. Модель получает проверенные признаки клиента и описание допустимых действий и возвращает несколько последовательностей в структурированном формате. Содержательная связность последовательности (например, сначала ознакомление, затем предложение) может быть сильной стороной LLM, но кандидат попадает в сравнение только после проверки условий задачи, а выбор выполняется по~\eqref{eq:plan}. Такая схема ограничивает риск концентрации на одном варианте, отмеченный в~\cite{hofer2026}.

\subsection{Генерация и проверка текста}

Бандит обучается на метках, а не на конкретных формулировках, что сокращает пространство действий и следует разделению ролей из~\cite{song2025}. LLM получает выбранные метки, описание предложения и подтверждённые факты о клиенте, а возвращает заголовок и текст. Затем выполняются проверки: длина, запрещённые темы, наличие обязательных условий предложения и соответствие меткам. Каждое фактическое утверждение проверяется по входным данным. По аналогии с~\cite{fostiropoulos2026} вычисляется доля подтверждённых утверждений
\begin{equation}
  G = \frac{|\text{подтверждённые утверждения}|}{|\text{все утверждения}|}.
  \label{eq:grounded}
\end{equation}
При $G < 1$ или непройденной проверке текст генерируется повторно, а затем заменяется заранее подготовленным шаблоном той же метки.

\subsection{Данные и полусинтетическая среда}

Открытых данных, где одновременно есть несколько сервисов, каналы, время, частота и контрольная группа без контактов, найти не удалось. Поэтому используется полусинтетическая среда на основе Open Bandit Dataset~\cite{saito2021}. Этот набор собран на платформе электронной коммерции и содержит признаки пользователей, показанные предметы, клики и вероятности показа для политик Random и Bernoulli TS. Предметы рассматриваются как предложения, а кампании набора условно как сервисы.

По логам политики Random обучается модель клика $q(x, p)$. Истинный отклик задаётся так:
\begin{align}
  \eta^{*}(x, \varnothing) &= \mathrm{logit}\, b(x), \label{eq:simbase} \\
  \eta^{*}(x, a) &= \mathrm{logit}\, b(x) + \beta_0 + \alpha\,\psi(x, p) + \beta_c + \beta_s(x) \nonumber \\
  &\quad + \beta_{\ell}(x) + \beta_{c\ell} - \rho\, f(h), \label{eq:simact}
\end{align}
где базовая вероятность $b(x) = r\,\bar{q}(x)$ получена масштабированием средней вероятности клика по предложениям, $\psi(x, p) = \mathrm{logit}\, q(x, p) - \mathrm{logit}\, \bar{q}(x)$. Эффекты канала, слота и меток, включая неоднородность по клиентам и их парное взаимодействие $\beta_{c\ell}$, порождаются симулятором. Функция утомления $f(h)$ растёт с числом недавних контактов и убывает со временем. Таким образом, реальные логи определяют только базовый уровень и предпочтения по предложениям, а эффект контакта является допущением.

Для снижения зависимости выводов от этих допущений проводится анализ по сетке значений $r$, $\beta_0$, $\rho$, размера неоднородности и горизонта $T$. Истинный эффект известен, что позволяет проверить не только награду, но и точность его оценки.

\subsection{Методика оценки}

\textit{Базовые стратегии}: случайная политика с тем же лимитом контактов; расписание по фиксированному правилу; жадная оптимизация отклика без молчания; бандит с независимой рукой на каждое сочетание; LinUCB~\cite{li2010} по признакам сочетания. Вклад компонентов оценивается поочерёдным отключением молчания, взаимодействий, признаков утомления, планирования и кандидатов от LLM.

\textit{Метрики в среде}: сожаление~\eqref{eq:regret}, истинная накопленная награда, доля истинного эффекта от оракула, число контактов на клиента, число нарушений ограничений. Точность оценки эффекта измеряется среднеквадратичной ошибкой $\tau_{\mu}$ относительно истинного $\tau$, отдельно для каждого модуля. Для каждой конфигурации выполняется 20 запусков; доверительные интервалы строятся бутстрэпом.

\textit{Проверка на логах.} Модуль выбора предложения дополнительно оценивается на реальных логах методом воспроизведения~\cite{li2011} и взвешиванием по обратной вероятности
\begin{equation}
  \hat{V}_{\mathrm{IPS}}(\pi) = \frac{1}{n} \sum_{i=1}^{n} \frac{\pi(p_i \mid x_i)}{\pi_0(p_i \mid x_i)}\, y_i,
  \label{eq:ips}
\end{equation}
где $\pi_0$~--- политика сбора данных, а также двойной устойчивой оценкой~\cite{dudik2011}. Этот шаг оценивает клики, а не инкрементальный эффект, и служит только проверкой реалистичности.

\textit{Тексты} оцениваются на выборке решений по доле~\eqref{eq:grounded} до и после проверок, соответствию меткам и числу обращений к LLM. Сравниваются шаблоны, LLM без ограничений и LLM с проверками. Влияние текста на поведение реальных клиентов в данной постановке не определяется.

Выводы работы будут ограничены полусинтетической средой: численный прирост нельзя перенести на реальный продукт. Для такого переноса потребуется полевой эксперимент с рандомизированной контрольной группой.

% ============================================================
\begin{thebibliography}{99}

\bibitem{abboud2025} S. Abboud, E. Hanna, O. Jeunen, V. Raheja, and S. Wheeler, “Agentic personalisation of cross-channel marketing experiences,” arXiv:2506.16429, 2025.

\bibitem{sato2019} M. Sato, J. Singh, S. Takemori, T. Sonoda, Q. Zhang, and T. Ohkuma, “Uplift-based evaluation and optimization of recommenders,” in \textit{Proc. 13th ACM Conference on Recommender Systems}, 2019, pp.~296--304, doi: 10.1145/3298689.3347018.

\bibitem{yancey2020} K. P. Yancey and B. Settles, “A sleeping, recovering bandit algorithm for optimizing recurring notifications,” in \textit{Proc. 26th ACM SIGKDD International Conference on Knowledge Discovery \& Data Mining}, 2020, pp.~3008--3016, doi: 10.1145/3394486.3403351.

\bibitem{fostiropoulos2026} I. Fostiropoulos et al., “GISTBench: Evaluating LLM user understanding via evidence-based interest verification,” arXiv:2603.29112, 2026.

\bibitem{zhang2026} L. Zhang, J. Liu, and T. August, “Re-centering humans in LLM personalization,” arXiv:2606.06614, 2026.

\bibitem{li2010} L. Li, W. Chu, J. Langford, and R. E. Schapire, “A contextual-bandit approach to personalized news article recommendation,” in \textit{Proc. 19th International Conference on World Wide Web}, 2010, pp.~661--670, doi: 10.1145/1772690.1772758.

\bibitem{li2011} L. Li, W. Chu, J. Langford, and X. Wang, “Unbiased offline evaluation of contextual-bandit-based news article recommendation algorithms,” in \textit{Proc. 4th ACM International Conference on Web Search and Data Mining}, 2011, pp.~297--306, doi: 10.1145/1935826.1935878.

\bibitem{thompson1933} W. R. Thompson, “On the likelihood that one unknown probability exceeds another in view of the evidence of two samples,” \textit{Biometrika}, vol.~25, no.~3--4, pp.~285--294, 1933, doi: 10.1093/biomet/25.3-4.285.

\bibitem{agrawal2013} S. Agrawal and N. Goyal, “Thompson sampling for contextual bandits with linear payoffs,” in \textit{Proc. 30th International Conference on Machine Learning}, PMLR, vol.~28, no.~3, 2013, pp.~127--135.

\bibitem{russo2018} D. J. Russo, B. Van Roy, A. Kazerouni, I. Osband, and Z. Wen, “A tutorial on Thompson sampling,” \textit{Foundations and Trends in Machine Learning}, vol.~11, no.~1, pp.~1--96, 2018, doi: 10.1561/2200000070.

\bibitem{hill2017} D. N. Hill, H. Nassif, Y. Liu, A. Iyer, and S. V. N. Vishwanathan, “An efficient bandit algorithm for realtime multivariate optimization,” in \textit{Proc. 23rd ACM SIGKDD International Conference on Knowledge Discovery and Data Mining}, 2017, pp.~1813--1821, doi: 10.1145/3097983.3098184.

\bibitem{mcinerney2018} J. McInerney et al., “Explore, exploit, and explain: Personalizing explainable recommendations with bandits,” in \textit{Proc. 12th ACM Conference on Recommender Systems}, 2018, pp.~31--39, doi: 10.1145/3240323.3240354.

\bibitem{kunzel2019} S. R. Künzel, J. S. Sekhon, P. J. Bickel, and B. Yu, “Metalearners for estimating heterogeneous treatment effects using machine learning,” \textit{Proceedings of the National Academy of Sciences}, vol.~116, no.~10, pp.~4156--4165, 2019, doi: 10.1073/pnas.1804597116.

\bibitem{dudik2011} M. Dudík, J. Langford, and L. Li, “Doubly robust policy evaluation and learning,” in \textit{Proc. 28th International Conference on Machine Learning}, 2011, pp.~1097--1104, arXiv:1103.4601.

\bibitem{saito2021} Y. Saito, S. Aihara, M. Matsutani, and Y. Narita, “Open Bandit Dataset and Pipeline: Towards realistic and reproducible off-policy evaluation,” in \textit{NeurIPS Datasets and Benchmarks Track}, 2021, arXiv:2008.07146.

\bibitem{matz2024} S. C. Matz, J. D. Teeny, S. S. Vaid, H. Peters, G. M. Harari, and M. Cerf, “The potential of generative AI for personalized persuasion at scale,” \textit{Scientific Reports}, vol.~14, art.~4692, 2024, doi: 10.1038/s41598-024-53755-0.

\bibitem{song2025} H. Song et al., “Tailored behavior-change messaging for physical activity: Integrating contextual bandits and large language models,” arXiv:2506.07275, 2025.

\bibitem{hofer2026} D. P. Hofer et al., “Structured exploration vs. generative flexibility: A field study comparing bandit and LLM architectures for personalised health behaviour interventions,” arXiv:2603.06330, 2026.

\bibitem{sosedka2024} А. В. Соседка, П. И. Ибрагимов, «Количественный и качественный анализ аудитории Telegram в разрезе рекомендаций с использованием больших языковых моделей», в сб. \textit{Искусственный интеллект в промышленных, коммерческих, медицинских и финансовых приложениях}, вып.~2. М.: МИСИС, 2024, с.~116--120.

\bibitem{chetvertkov2026} Е. С. Четвертков, «Сравнительный анализ мультимодальных архитектур в задаче текстового поиска 3D-моделей LEGO», в сб. \textit{Искусственный интеллект в промышленных, коммерческих, медицинских и финансовых приложениях}, вып.~6, под ред. А. Р. Ефимова. М.: НИТУ «МИСИС», 2026, с.~141--146.

\bibitem{zaitsev2026} М. В. Зайцев, Д. В. Чун, «Исследование методов восстановления трехмерной сцены из одного изображения», в сб. \textit{Искусственный интеллект в промышленных, коммерческих, медицинских и финансовых приложениях}, вып.~6, под ред. А. Р. Ефимова. М.: НИТУ «МИСИС», 2026, с.~51--59.

\bibitem{artemyeva2022} Л. А. Артемьева, А. А. Дряженков, М. М. Потапов, «Устойчивое решение задачи квадратичной минимизации с неравномерно возмущённым оператором регуляризованным градиентным методом», \textit{Журнал вычислительной математики и математической физики}, т.~62, №~1, с.~12--22, 2022, doi: 10.1134/S0965542522010031.

\end{thebibliography}

\end{document}
